\section{Technical Implementation Details}
\label{sec:impl-details}

This appendix elaborates on the core algorithmic and implementation details of the LDM framework, covering Gaussian process surrogate modelling, acquisition function design and batch sampling schemes. All implementations strictly follow the acquisition-tilted search framework presented in the main text, and only expand on the specific execution logic for surrogate model training and sampling.

\subsection{Gaussian Process Surrogate}
\label{sec:llm-bo}

LDM employs a Gaussian Process (GP) as the probabilistic surrogate model. It fits the posterior distribution of the black-box reward function from historical experimental observations, and outputs calibrated predictive means and epistemic uncertainty to provide quantitative inputs for acquisition function computation.

\subsubsection{Prior Specification}
Let $\mathcal{X}$ denote the design space, $R(x)$ the true black-box reward function. Each experimental observation contains independent Gaussian noise:
\[
r_i = R(x_i) + \epsilon_i,\quad \epsilon_i \sim \mathcal{N}(0, \sigma_n^2)
\]
where $\sigma_n^2$ is the observation noise variance.

The GP prior is defined as:
\[
R(x) \sim \mathcal{GP}\left(m(x),\, k(x, x')\right)
\]
where:
\begin{itemize}
  \item The prior mean function $m(x)$ adopts a constant prior, i.e. $m(x) = m_0$, with $m_0$ a scalar parameter that can be optimised from data;
  \item $k(x, x')$ is a positive-definite kernel function that characterises the correlation between samples in the design space. Its specific form is adapted to the design space, for example an RBF kernel for continuous parameter spaces or a string kernel for sequence spaces;
  \item The kernel hyperparameters, observation noise variance $\sigma_n^2$ and prior constant $m_0$ together form the set of parameters to be estimated for the GP.
\end{itemize}

\subsubsection{Posterior Inference}
Let $\mathcal{D}_t = \{(x_i, r_i)\}_{i=1}^{t}$ denote the cumulative historical observation dataset at iteration $t$. We define the $t \times t$ kernel matrix $\mathbf{K}_t$ and the $t$-dimensional cross-covariance vector $\mathbf{k}_t(x)$:
\[
\mathbf{K}_t =
\begin{bmatrix}
k(x_1, x_1) & k(x_1, x_2) & \dots & k(x_1, x_t) \\
k(x_2, x_1) & k(x_2, x_2) & \dots & k(x_2, x_t) \\
\vdots      & \vdots      & \ddots & \vdots      \\
k(x_t, x_1) & k(x_t, x_2) & \dots & k(x_t, x_t)
\end{bmatrix},\quad
\mathbf{k}_t(x) =
\begin{bmatrix}
k(x_1, x) \\
k(x_2, x) \\
\vdots \\
k(x_t, x)
\end{bmatrix} \in \mathbb{R}^t
\]

By Bayesian conditioning, the posterior reward distribution for any candidate $x \in \mathcal{X}$ is Gaussian, with closed-form expressions for the posterior predictive mean and predictive variance:
\begin{align}
\mu_t(x) &= m_0 + \mathbf{k}_t(x)^\top \left(\mathbf{K}_t + \sigma_n^2 \mathbf{I}\right)^{-1} \left(\mathbf{r}_t - m_0 \cdot \mathbf{1}\right) \label{eq:postmean} \\
\sigma_t^2(x) &= k(x, x) - \mathbf{k}_t(x)^\top \left(\mathbf{K}_t + \sigma_n^2 \mathbf{I}\right)^{-1} \mathbf{k}_t(x). \label{eq:postvar}
\end{align}
where $\mathbf{r}_t = (r_1, \dots, r_t)^\top$ is the vector of observed rewards, and $\mathbf{1}$ is an all-ones vector of length $t$. The posterior mean estimates the expected reward of a candidate, while the posterior variance quantifies the epistemic uncertainty of the prediction.

\subsubsection{Hyperparameter Estimation}
The GP hyperparameters are fitted via \emph{Type-II maximum likelihood estimation}, whereby optimal hyperparameter values are determined by maximising the marginal log-likelihood of the observed data. The marginal log-likelihood of the observed data $\mathcal{D}_t$ is given by:
\[
\log p(\mathbf{r}_t \mid \mathcal{D}_t) = -\frac{1}{2} (\mathbf{r}_t - m_0 \cdot \mathbf{1})^\top \left(\mathbf{K}_t + \sigma_n^2 \mathbf{I}\right)^{-1} (\mathbf{r}_t - m_0 \cdot \mathbf{1}) - \frac{1}{2}\log\left|\mathbf{K}_t + \sigma_n^2 \mathbf{I}\right| - \frac{t}{2}\log 2\pi
\]

Optimisation is performed using gradient-based algorithms such as L-BFGS. The parameters to be optimised include kernel hyperparameters such as lengthscales and output variance, observation noise variance $\sigma_n^2$, and the prior constant $m_0$. This method automatically adapts to the smoothness and noise level of the design space, improving the fitting accuracy and uncertainty calibration of the surrogate model. Hyperparameter optimisation is re-run after each iteration with new observations to update the surrogate model.

\subsection{Acquisition Functions}
Acquisition functions quantify the experimental value of each candidate based on the GP posterior mean and uncertainty, and act as the core value signal guiding search direction and balancing exploitation and exploration. LDM supports three acquisition functions for single- and multi-objective settings, all of which can be directly embedded into the acquisition-tilted search framework.

\subsubsection{Expected Improvement}
Expected Improvement (EI) \citep{jonesEfficientGlobalOptimization1998} measures the expected gain of a candidate relative to the current best observed value. It is one of the most widely used acquisition functions in single-objective optimisation, and takes the form:
\begin{equation}
\alpha^{\text{EI}}(x) = \left(\mu_t(x) - r_t^* - \xi\right) \Phi(z) + \sigma_t(x) \varphi(z), \quad z = \frac{\mu_t(x) - r_t^* - \xi}{\sigma_t(x)}
\label{eq:ei}
\end{equation}
where $r_t^* = \max_{i \le t} r_i$ is the best reward observed so far, $\xi \ge 0$ is an exploration-adjustment hyperparameter, and $\Phi(\cdot)$ and $\varphi(\cdot)$ denote the cumulative distribution function and probability density function of the standard normal distribution, respectively. EI naturally balances exploitation and exploration, prioritising candidates with higher expected returns.

\subsubsection{Upper Confidence Bound}
The Upper Confidence Bound (UCB) \citep{srinivas2010gaussian} computes a weighted sum of the predictive mean and uncertainty. It has a concise form and enjoys rigorous theoretical guarantees on cumulative regret, and takes the form:
\begin{equation}
\alpha^{\text{UCB}}(x) = \mu_t(x) + \sqrt{\beta_t}\, \sigma_t(x)
\label{eq:ucb}
\end{equation}
where $\beta_t > 0$ is the exploration weight coefficient, which may be set via theoretical formulae or tuned as a hyperparameter to control the exploration intensity of the algorithm. By adjusting $\beta_t$, UCB can flexibly switch between exploitation-biased and exploration-biased search strategies.

\subsubsection{Expected Hypervolume Improvement}

For multi-objective optimisation tasks, \textbf{Expected Hypervolume Improvement (EHVI)} \citep{ponweiserMultiobjectiveOptimizationLimited2008} is adopted as the acquisition function. EHVI computes the expected hypervolume gain of the current empirical Pareto front after adding a candidate, and automatically balances exploitation of favourable objective trade-offs and exploration of under-sampled design regions, without requiring pre-specified objective weights.

In the multi-objective setting, the core acquisition-tilted search framework remains unchanged; only the scalar acquisition value $a_t(x)$ is replaced with the EHVI value, enabling seamless extension of LDM to multi-objective discovery.

\subsection{Batch Acquisition Sampling}
\label{sec:batch}

In parallel experimental settings, multiple candidates must be selected for simultaneous evaluation at each iteration. LDM adopts Gumbel-top-$k$ sampling to draw batches of candidates directly from the acquisition-tilted target distribution $\pi_t(x)$, without requiring additional diversity regularisation terms.

The implementation procedure is as follows:
\begin{enumerate}
  \item For each candidate $x_i$ in the pool generated by the LLM, compute its unnormalised tilted weight:
  \[
  w_i = p_{\theta,\alpha}(x_i \mid \mathcal{C}_t) \cdot \exp\left\{\eta \cdot a_t(x_i)\right\}
  \]
  where $p_{\theta,\alpha}(x_i \mid \mathcal{C}_t)$ is the conditional generation probability of the LLM, $a_t(x_i)$ is the acquisition function value, and $\eta$ is the acquisition tilt strength coefficient.

  \item Sample independent Gumbel noise $g_i \sim \text{Gumbel}(0, 1)$ for each weight, rank candidates by $\log w_i + g_i$ in descending order, and select the top $k$ candidates to form the evaluation batch. Each selected candidate is removed from the pool after selection.
\end{enumerate}

Batch samples produced by this method strictly follow the target tilted distribution, naturally preserving structural diversity while maintaining acquisition value, and are compatible with batch evaluation pipelines in parallel experiments.